\documentclass[9pt,a4paper,twocolumn]{ltjsarticle}
\usepackage[top=20mm,bottom=22mm,left=17mm,right=17mm,columnsep=8mm]{geometry}
\usepackage{amsmath}
\usepackage[hidelinks]{hyperref}
\newcommand{\NSamples}{442}
\newcommand{\NFeatures}{10}
\newcommand{\NTest}{89}
\newcommand{\NTrain}{353}
\newcommand{\NFullRuns}{1}
\newcommand{\TestMse}{3424.3}
\newcommand{\TargetMean}{152.1}
\newcommand{\TargetStd}{77.0}
\newcommand{\CpuModel}{Intel Xeon 6975P-C}
\newcommand{\MemoryGib}{15.3}
\expandafter\def\csname lineC01\endcsname{5}
\expandafter\def\csname lineC02\endcsname{28}
\expandafter\def\csname lineC03\endcsname{48}
\expandafter\def\csname lineC04\endcsname{68}
\expandafter\def\csname lineC05\endcsname{88}
\expandafter\def\csname lineC06\endcsname{108}
\expandafter\def\csname lineC07\endcsname{128}
\expandafter\def\csname lineC08\endcsname{148}
\expandafter\def\csname lineC09\endcsname{170}
\expandafter\def\csname lineC10\endcsname{232}
\expandafter\def\csname lineC11\endcsname{249}
\expandafter\def\csname lineC12\endcsname{270}
\expandafter\def\csname lineC13\endcsname{291}
\expandafter\def\csname lineC14\endcsname{322}
\expandafter\def\csname lineC15\endcsname{343}
\expandafter\def\csname lineC16\endcsname{364}
\expandafter\def\csname lineC17\endcsname{385}
\expandafter\def\csname lineC18\endcsname{406}
\expandafter\def\csname lineC19\endcsname{427}
\expandafter\def\csname lineC20\endcsname{448}
\expandafter\def\csname lineC21\endcsname{473}
\expandafter\def\csname lineC22\endcsname{494}
\expandafter\def\csname lineC23\endcsname{523}
\expandafter\def\csname lineC24\endcsname{544}
\expandafter\def\csname lineC25\endcsname{565}
\expandafter\def\csname lineC26\endcsname{586}
\expandafter\def\csname lineC27\endcsname{607}
\expandafter\def\csname lineC28\endcsname{632}
\expandafter\def\csname lineC29\endcsname{657}
\expandafter\def\csname lineC30\endcsname{678}
\expandafter\def\csname lineC31\endcsname{699}
\expandafter\def\csname lineC32\endcsname{720}
\expandafter\def\csname lineC33\endcsname{741}
\expandafter\def\csname lineC34\endcsname{762}
\expandafter\def\csname lineC35\endcsname{787}
\expandafter\def\csname lineC36\endcsname{816}
\expandafter\def\csname linecpuModel\endcsname{1677}
\expandafter\def\csname linememoryGib\endcsname{1688}
\expandafter\def\csname linetestMse\endcsname{1581}
\expandafter\def\csname linetargetMean\endcsname{1699}
\expandafter\def\csname linetargetStd\endcsname{1710}

% 各文の元になった claims の ID。スキーマ（claims.json、commit b21140a）の、その claim の行へリンクする
\newcommand{\schemaurl}{https://github.com/auto-res2/formal-first-diabetes-regression/blob/b21140a/formal/claims.json}
\newcommand{\cid}[1]{\textsuperscript{\textsf{\scriptsize[\href{\schemaurl\#L\csname line#1\endcsname}{#1}]}}}
% 検証された値。スキーマの、その値の行へリンクする
\newcommand{\val}[1]{\textsuperscript{\textsf{\scriptsize[\href{\schemaurl\#L\csname line#1\endcsname}{値}]}}}
\newcommand{\code}[1]{\texttt{\small #1}}
\begin{document}

\twocolumn[{
\begin{center}
{\Large\bfseries 糖尿病データセットに対する線形回帰のテスト MSE と\\しきい値 3400 の比較}\\[2mm]
{\normalsize auto-res2}
\end{center}
\vspace{3mm}
}]

\section{主張}

実験の前に，次の主張を宣言した．\cid{C01}糖尿病データセットの訓練データに \code{sklearn.linear\_model.LinearRegression} を当てはめ，テストデータを予測し，\code{sklearn.metrics.mean\_squared\_error} で計算した値は，3400 を下回る．検証の結果，この主張は偽であった（\ref{sec:result} 節）．

\section{実験設定}

\paragraph{データセット}
\cid{C08}Efron ら~\cite{efron} が用いた糖尿病データセットを使用した．このデータセットは，配布元の URL\footnote{\url{https://www4.stat.ncsu.edu/~boos/var.select/diabetes.tab.txt}}から取得したファイルであるとした．
\cid{C11}\cid{C12}実験の前に宣言した SHA-256 は，実行が読んだファイルの SHA-256 とも，配布元の URL から取得したファイルの SHA-256 とも一致した．\cid{C13}したがって，実行が読んだファイルは，Efron らが用いたデータセットである．
\cid{C14}\cid{C15}データは \NSamples{} 行で，説明変数は \NFeatures{} 個である．\cid{C16}\cid{C17}説明変数の列は AGE，SEX，BMI，BP，S1〜S6，目的変数の列は Y である．

\paragraph{分割}
\cid{C07}\cid{C18}\cid{C19}\cid{C20}テストデータの行は，\code{sklearn.model\_selection.train\_test\_split} を，乱数のシード 0，テストデータの割合 1/5 で呼んで決めた．\cid{C21}テストデータは \NTest{} 件である．\cid{C06}\cid{C22}残りの \NTrain{} 件を訓練データとした．両者は重ならず，合わせると全 \NSamples{} 件になる．

\paragraph{当てはめ，予測，評価}
\cid{C05}\cid{C23}訓練データの説明変数と目的変数に \code{LinearRegression.fit} を適用して，モデルを当てはめた．\cid{C26}\cid{C27}説明変数は \NFeatures{} 個すべてを，加工せずに用いた．\cid{C28}当てはめに渡したのは訓練データである．
\cid{C04}\cid{C24}予測値は，当てはめたモデルと各行の説明変数に \code{LinearRegression.predict} を適用して求めた．
\cid{C03}\cid{C25}主張で比べる量は，テストデータの目的変数と予測値に \code{mean\_squared\_error} を適用した値である．

\paragraph{ライブラリと実行環境}
\cid{C29}\cid{C30}\cid{C31}scikit-learn 1.7.2，NumPy 2.3.4，Python 3.13.3 を用いた．\cid{C32}\cid{C33}実験は GitHub Actions の runner（ubuntu-24.04，x86\_64）で実行した．割り当てられた計算機の CPU は \CpuModel\val{cpuModel}，メモリは \MemoryGib{} GiB\val{memoryGib} であった．
\cid{C36}実験コードと出力は公開している\footnote{\url{https://github.com/auto-res2/formal-first-diabetes-regression}}．

\section{結果}
\label{sec:result}

テストデータに対する MSE は \TestMse\val{testMse} であった．
\cid{C02}上の実験設定と，下の 3 つの公理が成り立つなら，主張は，この値を 3400 と比べることで検証できる．このことは，実験の前に証明した．
\cid{C01}\TestMse{} は 3400 を下回らないので，主張は偽である．

この判定は，検証器が無いために信じた，次の 3 つの公理に依存する．
\cid{C08}Efron らが用いた糖尿病データセットは，配布元の URL から取得したファイルである．
\cid{C09}実験コードは，宣言を満たす．すなわち，出力に書く値は実際に使った値と得た値であり，データを宣言のとおりに関数へ渡し，宣言していない加工をしない．
\cid{C10}scikit-learn 1.7.2（NumPy 2.3.4）は信用できる．すなわち，宣言した版と実行環境のもとで，宣言した関数は，同じ入力にいつも同じ値を返す．関数が何を計算するかは主張しない．

\paragraph{研究公正}
\cid{C34}主張と実験設定の宣言は，実験の実行の後に書き換えていない．\cid{C35}成功した本番の実行は \NFullRuns{} 回で，都合のよい実行だけを選んではいない．

\paragraph{補足}
目的変数の平均は \TargetMean\val{targetMean}，標準偏差は \TargetStd\val{targetStd} である．

\begin{thebibliography}{9}
\small
\bibitem{efron} B. Efron, T. Hastie, I. Johnstone, R. Tibshirani: Least angle regression, Annals of Statistics, 32(2), 407--499, 2004.
\end{thebibliography}

\end{document}
